% GENERATED FILE. Edit canonical lessons, then run npm run generate:books.
% canonical-source-hash: fnv1a64:2601ad1a40e8ef81
% canonical-lessons: FA-C526-mojassame-saz, FA-C526-memari, FA-C526-miniyator, FA-C526-tarrahi-kardan, FA-C526-khalq-kardan

\chapter{Art and Design}
\label{ch:fa-art-and-design}
\hlchaptermodality{526}

\begin{chapteropening}
\textbf{By the end of this chapter:} \emph{I can talk about art, architecture and design.}

Complete the last lesson of chapter 526: I can talk about art, architecture and design.
\end{chapteropening}

\section[\texorpdfstring{\fa{مجسمه} \fa{ساز} (mojassame sâz)}{مجسمه ساز (mojassame sâz)}]{\fa{مجسمه} \fa{ساز} (mojassame sâz) — a sculptor}

\label{lesson:FA-C526-mojassame-saz}



\begin{quote}
Before the new one: say the Persian for watercolour, then the Persian for a gallery.
\end{quote}

\subsection*{You'll want to know: \fa{مجسمه} \fa{ساز}}
\textbf{\fa{مجسمه} \fa{ساز}} (\emph{mojassame sâz}) — \textquotedblleft{}a sculptor\textquotedblright{}.

\begin{quote}
\fa{مجسمه} \fa{ساز} \fa{با} \fa{سنگ} \fa{کار} \fa{میکند}.
\end{quote}

\begin{quote}
\emph{mojassame sâz bâ sang kâr mikonad.}
\end{quote}

\begin{quote}
The sculptor works in stone.
\end{quote}

\emph{mojassame} 'statue' + \emph{sâz} 'maker' (stem of \emph{sâkhtan}).

\begin{cousinweb}
\emph{mojassame sâz} is built in Persian on \emph{mojassame}, 'statue', which was borrowed from Arabic \emph{mujassama} (EtymDB), and \emph{sâz}, the present stem of native \emph{sâkhtan}, 'to make'.
\end{cousinweb}

\begin{grammarlens}[title={What you've built}]
One more word for telling, explaining and describing at B1.
\end{grammarlens}

\subsection*{Your turn}
\begin{itemize}
\raggedright
  \item \emph{Say it:} \emph{mojassame sâz}
  \item \emph{Say it:} \emph{mojassame sâz}, once more
  \item \emph{Say it:} a sentence of your own with \emph{mojassame sâz}
  \item \emph{Recall:} say the Persian for watercolour, then the Persian for a gallery, then say \emph{mojassame sâz} again
\end{itemize}

\begin{tcolorbox}[breakable,colback=teal!4,colframe=teal!35!black,title={Before you move on}]
What does \fa{مجسمه} \fa{ساز} mean? (\textquotedblleft{}A sculptor\textquotedblright{}.) Say it once more.
\end{tcolorbox}
\section[\texorpdfstring{\fa{معماری} (me'mâri)}{معماری (me'mâri)}]{\fa{معماری} (me'mâri) — architecture}

\label{lesson:FA-C526-memari}



\begin{quote}
Before the new one: say the Persian for a gallery, then the Persian for a sculptor.
\end{quote}

\subsection*{You'll want to know: \fa{معماری}}
\textbf{\fa{معماری}} (\emph{me'mâri}) — \textquotedblleft{}architecture\textquotedblright{}.

\begin{quote}
\fa{او} \fa{در} \fa{دانشگاه} \fa{معماری} \fa{میخواند}.
\end{quote}

\begin{quote}
\emph{u dar dâneshgâh me'mâri mikhânad.}
\end{quote}

\begin{quote}
She studies architecture at university.
\end{quote}

An architect is \emph{me'mâr}.

\begin{cousinweb}
\emph{me'mâri} comes from Arabic \emph{miʿmārī} (EtymDB), built on \emph{miʿmār}, 'builder', which Persian keeps as \emph{me'mâr}, 'architect'.
\end{cousinweb}

\begin{grammarlens}[title={What you've built}]
One more word for telling, explaining and describing at B1.
\end{grammarlens}

\subsection*{Your turn}
\begin{itemize}
\raggedright
  \item \emph{Say it:} \emph{me'mâri}
  \item \emph{Say it:} \emph{me'mâri}, once more
  \item \emph{Say it:} a sentence of your own with \emph{me'mâri}
  \item \emph{Recall:} say the Persian for a gallery, then the Persian for a sculptor, then say \emph{me'mâri} again
\end{itemize}

\begin{tcolorbox}[breakable,colback=teal!4,colframe=teal!35!black,title={Before you move on}]
What does \fa{معماری} mean? (\textquotedblleft{}Architecture\textquotedblright{}.) Say it once more.
\end{tcolorbox}
\section[\texorpdfstring{\fa{مینیاتور} (miniyâtor)}{مینیاتور (miniyâtor)}]{\fa{مینیاتور} (miniyâtor) — a miniature painting}

\label{lesson:FA-C526-miniyator}



\begin{quote}
Before the new one: say the Persian for a sculptor, then the Persian for architecture.
\end{quote}

\subsection*{You'll want to know: \fa{مینیاتور}}
\textbf{\fa{مینیاتور}} (\emph{miniyâtor}) — \textquotedblleft{}a miniature painting\textquotedblright{}.

\begin{quote}
\fa{در} \fa{موزه} \fa{مینیاتورهای} \fa{قدیمی} \fa{دیدیم}.
\end{quote}

\begin{quote}
\emph{dar muze miniyâtorhâ-ye qadimi didim.}
\end{quote}

\begin{quote}
We saw old miniatures in the museum.
\end{quote}

Persian miniature painting is an old art of book illustration.

\begin{cousinweb}
\emph{miniyâtor} is a modern loan, usually traced to French \emph{miniature}. EtymDB has no entry for the Persian word.
\end{cousinweb}

\begin{grammarlens}[title={What you've built}]
One more word for telling, explaining and describing at B1.
\end{grammarlens}

\subsection*{Your turn}
\begin{itemize}
\raggedright
  \item \emph{Say it:} \emph{miniyâtor}
  \item \emph{Say it:} \emph{miniyâtor}, once more
  \item \emph{Say it:} a sentence of your own with \emph{miniyâtor}
  \item \emph{Recall:} say the Persian for a sculptor, then the Persian for architecture, then say \emph{miniyâtor} again
\end{itemize}

\begin{tcolorbox}[breakable,colback=teal!4,colframe=teal!35!black,title={Before you move on}]
What does \fa{مینیاتور} mean? (\textquotedblleft{}A miniature painting\textquotedblright{}.) Say it once more.
\end{tcolorbox}
\section[\texorpdfstring{\fa{طراحی} \fa{کردن} (tarrâhi kardan)}{طراحی کردن (tarrâhi kardan)}]{\fa{طراحی} \fa{کردن} (tarrâhi kardan) — to design}

\label{lesson:FA-C526-tarrahi-kardan}



\begin{quote}
Before the new one: say the Persian for architecture, then the Persian for a miniature painting.
\end{quote}

\subsection*{You'll want to know: \fa{طراحی} \fa{کردن}}
\textbf{\fa{طراحی} \fa{کردن}} (\emph{tarrâhi kardan}) — \textquotedblleft{}to design\textquotedblright{}.

\begin{quote}
\fa{او} \fa{لباسهای} \fa{زنانه} \fa{طراحی} \fa{میکند}.
\end{quote}

\begin{quote}
\emph{u lebâshâ-ye zanâne tarrâhi mikonad.}
\end{quote}

\begin{quote}
She designs women's clothes.
\end{quote}

\emph{tarrâhi} 'design' + \emph{kardan}; built on \emph{tarh} 'plan, design'.

\begin{cousinweb}
\emph{tarrâhi} is built in Persian on \emph{tarrâh}, 'designer', from \emph{tarh}, 'plan, design', a loan usually traced to Arabic \emph{ṭarḥ}. EtymDB has no entry for it. \emph{kardan} is inherited.
\end{cousinweb}

\begin{grammarlens}[title={What you've built}]
One more word for telling, explaining and describing at B1.
\end{grammarlens}

\subsection*{Your turn}
\begin{itemize}
\raggedright
  \item \emph{Say it:} \emph{tarrâhi kardan}
  \item \emph{Say it:} \emph{tarrâhi kardan}, once more
  \item \emph{Say it:} a sentence of your own with \emph{tarrâhi kardan}
  \item \emph{Recall:} say the Persian for architecture, then the Persian for a miniature painting, then say \emph{tarrâhi kardan} again
\end{itemize}

\begin{tcolorbox}[breakable,colback=teal!4,colframe=teal!35!black,title={Before you move on}]
What does \fa{طراحی} \fa{کردن} mean? (\textquotedblleft{}To design\textquotedblright{}.) Say it once more.
\end{tcolorbox}
\section[\texorpdfstring{\fa{خلق} \fa{کردن} (khalq kardan)}{خلق کردن (khalq kardan)}]{\fa{خلق} \fa{کردن} (khalq kardan) — to create}

\label{lesson:FA-C526-khalq-kardan}



\begin{quote}
Before the new one: say the Persian for a miniature painting, then the Persian for to design.
\end{quote}

\subsection*{You'll want to know: \fa{خلق} \fa{کردن}}
\textbf{\fa{خلق} \fa{کردن}} (\emph{khalq kardan}) — \textquotedblleft{}to create\textquotedblright{}.

\begin{quote}
\fa{او} \fa{با} \fa{چند} \fa{رنگ} \fa{ساده} \fa{یک} \fa{نقاشی} \fa{زیبا} \fa{خلق} \fa{کرد}.
\end{quote}

\begin{quote}
\emph{u bâ chand rang-e sâde yek naqqâshi-ye zibâ khalq kard.}
\end{quote}

\begin{quote}
With a few simple colours she created a beautiful painting.
\end{quote}

\emph{khalq} 'creation' + \emph{kardan}. Formal; in speech \emph{dorost kardan} or \emph{sâkhtan}.

\begin{cousinweb}
\emph{khalq} comes from Arabic \emph{ḫalq} (EtymDB), which EtymDB glosses 'people, creature': what has been created. \emph{kardan} is inherited Persian.
\end{cousinweb}

\begin{grammarlens}[title={What you've built}]
One more word for telling, explaining and describing at B1.
\end{grammarlens}

\subsection*{Your turn}
\begin{itemize}
\raggedright
  \item \emph{Say it:} \emph{khalq kardan}
  \item \emph{Say it:} \emph{khalq kardan}, once more
  \item \emph{Say it:} a sentence of your own with \emph{khalq kardan}
  \item \emph{Recall:} say the Persian for a miniature painting, then the Persian for to design, then say \emph{khalq kardan} again
\end{itemize}

\begin{tcolorbox}[breakable,colback=teal!4,colframe=teal!35!black,title={Before you move on}]
What does \fa{خلق} \fa{کردن} mean? (\textquotedblleft{}To create\textquotedblright{}.) Say it once more.
\end{tcolorbox}
